\section{Experiments}
\begin{table*}[t!]
	\caption{Test log likelihood (in nats) for UCI datasets and BSDS300, with error bars corresponding to two standard deviations. FFJORD$^\dagger$, NAF$^\dagger$, Block-NAF$^\dagger$, and SOS$^\dagger$ report error bars across repeated runs rather than across the test set. Superscript$^\star$ indicates results are taken from the existing literature. For validation results which can be used for comparison during model development, see \cref{tab:uci-results-validation} in \cref{appendix:experimental-details-tabular-density-estimation}.
	}
	\label{tab:uci-results}
        \vspace*{-0.1in}
	\begin{center}
		\begin{small}
			\begin{sc}
				\begin{tabular}{lccccc}
					\toprule
					Model & POWER & GAS & HEPMASS & MINIBOONE & BSDS300 \\ 
					\midrule
					FFJORD$^{\star\dagger}$ & $ 0.46 \pm 0.01$ & $ \hphantom{0}8.59 \pm 0.12$ & $ -14.92 \pm 0.08$ & $ -10.43 \pm 0.04$ & $ 157.40 \pm 0.19$ \\
					Glow & $ 0.42 \pm 0.01 $ & $ 12.24 \pm 0.03 $ & $ -16.99 \pm 0.02 $ & $ -10.55 \pm 0.45 $ & $ 156.95 \pm 0.28 $ \\
					Q-NSF (C) & $ 0.64 \pm 0.01 $ & $ 12.80 \pm 0.02 $ & $ -15.35 \pm 0.02 $ & $ \hphantom{0}{-9.35} \pm 0.44 $ & $ 157.65 \pm 0.28 $ \\
					RQ-NSF (C) & $ 0.64 \pm 0.01 $ & $ 13.09 \pm 0.02 $ & $ -14.75 \pm 0.03 $ & $ \hphantom{0}{-9.67} \pm 0.47 $ & $ 157.54 \pm 0.28 $  \\
					\midrule
					MAF & $ 0.45 \pm 0.01 $ & $ 12.35 \pm 0.02 $ & $ -17.03 \pm 0.02 $ & $ -10.92 \pm 0.46 $ & $ 156.95 \pm 0.28 $ \\
					Q-NSF (AR) & $ 0.66 \pm 0.01 $ & $ 12.91 \pm 0.02 $ & $ -14.67 \pm 0.03 $ & $ \hphantom{0}{-9.72} \pm 0.47 $ & $ 157.42 \pm 0.28 $  \\
					NAF$^{\star\dagger}$ & $0.62 \pm 0.01$ & $11.96 \pm 0.33$ & $-15.09 \pm 0.40$ & $\hphantom{0}{-8.86} \pm 0.15$ & $157.73 \pm 0.04$ \\
					Block-NAF$^{\star\dagger}$ & $0.61 \pm 0.01$ & $12.06 \pm 0.09$ & $-14.71 \pm 0.38$ & $\hphantom{0}{-8.95} \pm 0.07$ & $157.36 \pm 0.03$ \\
					SOS$^{\star\dagger}$ & $ 0.60 \pm 0.01 $ & $ 11.99 \pm 0.41 $ & $ -15.15 \pm 0.10 $ & $ \hphantom{0}{-8.90} \pm 0.11 $ & $ 157.48 \pm 0.41 $ \\
					RQ-NSF (AR) & $ 0.66 \pm 0.01 $ & $ 13.09 \pm 0.02 $ & $ -14.01 \pm 0.03 $ & $ \hphantom{0}{-9.22} \pm 0.48 $ & $ 157.31 \pm 0.28 $ \\
					\bottomrule
				\end{tabular}
			\end{sc}
		\end{small}
	\end{center}
	\vspace{-1.0em}
\end{table*}

In our experiments, the neural network NN which computes the parameters of the elementwise transformations is a residual network \citep{he2016deep} with pre-activation  residual blocks \citep{he2016preact}. 
For autoregressive transformations, the layers must be masked so as to preserve autoregressive structure, and so we use the ResMADE architecture outlined by \citet{nash2019aem}. 
Preliminary results indicated only minor differences in setting the tail bound $ B $ within the range $ [1, 5] $, and so we fix a value $ B = 3 $ across experiments, and find this to work robustly.
We also fix the number of bins $ K = 8 $ across our experiments, unless otherwise noted.
We implement all invertible linear transformations using the LU-decomposition, where the permutation matrix $ \bfP $ is fixed at the beginning of training, and the product $ \bfL \bfU $ is initialized to the identity.
For all non-image experiments, we define a flow `step' as the composition of an invertible linear transformation with either a coupling or autoregressive transform, and we use 10 steps per flow in all our experiments, unless otherwise noted.
All flows use a standard-normal noise distribution.
We use the Adam optimizer \citep{kingma:2014:adam}, and anneal the learning rate according to a cosine schedule \citep{loshchilov:2016:sgdr}. 
In some cases, we find applying dropout \citep{srivastava:2014:dropout} in the residual blocks beneficial for regularization.
Full experimental details are provided in \cref{appendix:experimental-details}. Code is \ifanonymize included in the supplementary material. \else available online at \url{https://github.com/bayesiains/nsf}.

\subsection{Density estimation of tabular data}
\label{sec:tabular-density-estimation}
We first evaluate our proposed flows using a selection of datasets from the UCI machine-learning repository \cite{dua:2017:uci} and BSDS300 collection of natural images \citep{martin:2001:bsds}. 
We follow the experimental setup and pre-processing of \citet{Papamakarios:2017:maf}, who make their data available online \citep{Papamakarios:2018:maf_datasets}.
We also update their MAF results using our codebase with ResMADE and invertible linear layers instead of permutations, providing a stronger baseline.
For comparison, we modify the quadratic splines of \citet{Muller:2018:nis} to match the rational-quadratic transforms, by defining them on the range $ [-B, B] $ instead of $ [0, 1] $, and adding linear tails, also matching the boundary derivatives as in the rational-quadratic case. 
We denote this model Q-NSF\@. 
Our results are shown in \cref{tab:uci-results}, where the mid-rule separates flows with one-pass inverse from autoregressive flows.
We also include validation results for comparison during model development in \cref{tab:uci-results-validation} in \cref{appendix:experimental-details-tabular-density-estimation}.

Both RQ-NSF (C) and RQ-NSF (AR) achieve state-of-the-art results for a normalizing flow on the Power, Gas, and Hepmass datasets, tied with Q-NSF (C) and Q-NSF (AR) on the Power dataset. 
Moreover, RQ-NSF (C) significantly outperforms both Glow and FFJORD, achieving scores competitive with the best autoregressive models. 
These results close the gap between autoregressive flows and flows based on coupling layers, and demonstrate that, in some cases, it may not be necessary to sacrifice one-pass sampling for density-estimation performance.


\subsection{Improving the variational autoencoder}
\label{sec:vae}

Next, we examine our proposed flows in the context of the variational autoencoder \citep[VAE,][]{Kingma:2014:vae, Rezende:2014:vae}, where they can act as both flexible prior and approximate posterior distributions. 
For our experiments, we use dynamically binarized versions of the MNIST dataset of handwritten digits \citep{lecun:1998:mnist}, and the EMNIST dataset variant featuring handwritten letters \citep{cohen2017emnist}.
We measure the capacity of our flows to improve over the commonly used baseline of a standard-normal prior and diagonal-normal approximate posterior, as well as over either coupling or autoregressive distributions with affine transformations.
Quantitative results are shown in \cref{tab:vae-results}, and image samples
in \cref{appendix-additional-experimental-results}.

All models improve significantly over the baseline, but perform very similarly otherwise, with most featuring overlapping error bars. 
Considering the disparity in density-estimation performance in the previous section, this is likely due to flows with affine transformations being sufficient to model the latent space for these datasets, with little scope for RQ-NSF flows to demonstrate their increased flexibility. 
Nevertheless, it is worthwhile to highlight that RQ-NSF (C) is the first class of model which can potentially match the flexibility of autoregressive models, and which requires no modification for use as either a prior or approximate posterior, due to its one-pass invertibility.

\begin{table*}[t!]
	\caption{Variational autoencoder test-set results (in nats) for the evidence lower bound (ELBO) and importance-weighted estimate of the log likelihood (computed as by \citet{burda:2016:iwae} using 1000 importance samples). Error bars correspond to two standard deviations.}
	\label{tab:vae-results}
        \vspace*{-0.1in}
	\begin{center}
		\begin{small}
			\begin{sc}
				\begin{tabular}{lcccc}
					\toprule
					& \multicolumn{2}{c}{MNIST} & \multicolumn{2}{c}{EMNIST} \\ 
					\cmidrule(lr){2-3} \cmidrule(lr){4-5}
					Posterior/Prior & ELBO & $ \log p(\bfx) $ & ELBO & $ \log p(\bfx) $ \\
					\midrule
					Baseline & $ -85.61 \pm 0.51 $ & $ -81.31 \pm 0.43 $ & $ -125.89 \pm 0.41 $ & $ -120.88 \pm 0.38 $ \\
					\midrule
					Glow & $ -82.25 \pm 0.46 $ & $ -79.72 \pm 0.42 $ & $ -120.04 \pm 0.40 $ & $ -117.54 \pm 0.38 $ \\
					RQ-NSF (C) & $ -82.08 \pm 0.46 $ & $ -79.63 \pm 0.42 $ & $ -119.74 \pm 0.40 $ & $ -117.35 \pm 0.38 $ \\
					\midrule
					IAF/MAF & $ -82.56 \pm 0.48 $ & $ -79.95 \pm 0.43 $ & $ -119.85 \pm 0.40 $ & $ -117.47 \pm 0.38 $ \\ 
					RQ-NSF (AR) & $ -82.14 \pm 0.47 $ & $ -79.71 \pm 0.43 $ & $ -119.49 \pm 0.40 $ & $ -117.28 \pm 0.38 $ \\
					\bottomrule
				\end{tabular}
			\end{sc}
		\end{small}
	\end{center}
	\vspace{-1em}
\end{table*}

\subsection{Generative modeling of images}
\label{sec:gen_image_modeling}

Finally, we evaluate neural spline flows as generative models of images, measuring their capacity to improve upon baseline models with affine transforms.
In this section, we focus solely on flows with a one-pass inverse in the style of RealNVP \citep{Dinh:2017:rnvp} and Glow \citep{Kingma:2018:glow}.
We use the CIFAR-10 \citep{krizhevsky:2009:cifar10} and downsampled $64\times64$ ImageNet \citep{oord2016pixelrnn, russakovsky:2015:imagenet} datasets, with original 8-bit colour depth and with reduced 5-bit colour depth. 
We use Glow-like architectures with either affine (in the baseline model) or rational-quadratic coupling transforms, and provide full experimental detail in \cref{appendix:experimental-details}.
Quantitative results are shown in \cref{tab:img-results}, and samples are shown in \cref{fig:img-samples} and \cref{appendix-additional-experimental-results}.

RQ-NSF (C) improves upon the affine baseline in three out of four tasks, and the improvement is most significant on the 8-bit version of ImageNet64. At the same time, RQ-NSF (C) achieves scores that are competitive with the original Glow model, while significantly reducing the number of parameters required, in some cases by almost an order of magnitude.
\Cref{fig:img-samples} demonstrates that the model is capable of producing diverse, globally coherent samples which closely resemble real data. 
There is potential to further improve our results by replacing the uniform dequantization used in Glow with \textit{variational dequantization}, and using more powerful networks with gating and self-attention mechanisms to parameterize the coupling transforms, both of which are explored by \citet{Ho:2019:flowpp}.